\newif\ifglobal

%%%%%%%%%%%%%%%%%%%%%%%
%%% Compile to view %%%
%%%%%%%%%%%%%%%%%%%%%%%

%%%%%%%%%%%%%%%%%%%%%%%%%%%%%%%%%%%%%%%%%%%%%%%%%%%
%%% Readme:                                     %%%
%%% https://www.overleaf.com/read/bwdjvfjksvjy/ %%%
%%%%%%%%%%%%%%%%%%%%%%%%%%%%%%%%%%%%%%%%%%%%%%%%%%%

% >>> M?: marks a module setting number or important notice

% >>> M4: Set {./relative-path} to external files for this module <<<
\def \modulefiles {./33-I2S}
% To add an external file, use:
% (Replace '---' with actual filename)
% '\input{\modulefiles/tables/---.tex}' for tables
% '\includegraphics{\modulefiles/figures/---}' for figures
% '\subfile{\modulefiles/---}' for register files


% >>> M1: This module is in global project? <<<
% 'YES' - uncomment; 'NO' - comment out
\globaltrue

\ifglobal
    \documentclass[main\_\_EN.tex]{subfiles}
   
\else
    \documentclass[openany, 10pt]{book}

    % >>> M2: In ./00-shared/preamble-trm-module, also find
    %         settings for visibility of todo notes, watermark <<<
    \usepackage{preamble-trm-module}

    % >>> M3: Temporary labels in English,
    %         you can add your own labels there <<<
    \myexternaldocument{./00-shared/temp-labels-en}

\fi %\ifglobal


\setcounter{secnumdepth}{4}


% >>> M5: If this module needs any updates in
% './00-shared/preamble-trm-module.sty'
% (include additional package, etc.), contact your module PM
% or create the file './00-shared/changelog.tex' make the
% necessary updates yourself and document them in this file <<<


\begin{document}


% Sets language into which pre-defined bits of text are translated
\selectlanguage{english}

% Do not delete this block
\ifglobal
% Add the following front matter if
% this module is outside of global project
\else
    \InputIfFileExists{readme}{}{}
    {\let\clearpage\relax \listoftodos}
    {\let\clearpage\relax \listofdonetodos}
    \tableofcontents
    \listoftables
    \listoffigures
\fi


%%%%%%%%%%%%%%%%%%%%%%%%%
%%% TRM Module Begins %%%
%%%%%%%%%%%%%%%%%%%%%%%%%

\hypertarget{i2s}{}
\chapter{I2S Controller (I2S)}
\label{mod:i2s}


\section{Overview}
\chipname{} has a built-in I2S interface, which provides a flexible communication interface for streaming digital data in multimedia applications, especially digital audio applications.

The I2S standard bus defines three signals: a bit clock signal (BCK), a channel/word select signal (WS), and a serial data signal (SD). A basic I2S data bus has one master and one slave. The roles remain unchanged throughout the communication. The I2S module on \chipname{} provides separate transmit (TX) and receive (RX) units for high performance.

\section{Terminology}

To better illustrate the functionality of I2S, the following terms are used in this chapter. 

\begin{longtable}[c]{ R{4cm} p{10cm} }
\textbf{Master mode}     & As a master, I2S drives BCK/WS signals, and sends data to or receives data from a slave. \\ 
\textbf{Slave mode}      & As a slave, I2S is driven by BCK/WS signals, and receives data from or sends data to a master. \\
\textbf{Full-duplex} & There are two separate data lines. Transmitted and received data are carried simultaneously.\\
\textbf{Half-duplex} & Only one side, the master or the slave, sends data first, and the other side receives data. Sending data and receiving data can not happen at the same time.\\
\textbf{A-law and $\mu$-law} & A-law and \textrm{$\mu$}-law are compression/decompression algorithms in digital pulse code modulated (PCM) non-uniform quantization, which can effectively improve the signal-to-quantization noise ratio. \\
\textbf{TDM RX mode} & In this mode, pulse code modulated (PCM) data is received and stored into memory via direct memory access (DMA), utilizing time division multiplexing (TDM). The signal lines include: BCK, WS, and SD. Data from 16 channels at most can be received. TDM Philips standard, TDM MSB alignment standard, and TDM PCM standard are supported in this mode, depending on user configuration.\\
\textbf{Normal PDM RX mode} & In this mode, pulse density modulation (PDM) data is received and stored into memory via DMA. Used signals: WS and DATA. PDM standard is supported in this mode by user configuration.\\
\textbf{TDM TX mode} & In this mode, pulse code modulated (PCM) data is sent from memory via DMA, in a way of time division multiplexing (TDM). The signal lines include: BCK, WS, and DATA. Data up to 16 channels can be sent. TDM Philips standard, TDM MSB alignment standard, and TDM PCM standard are supported in this mode, depending on user configuration.\\
\textbf{Normal PDM TX mode} & In this mode, pulse density modulation (PDM) data is sent from memory via DMA. The signal lines include: WS and DATA. PDM standard is supported in this mode by user configuration.\\
\textbf{PCM-to-PDM TX mode} & In this mode, I2S as a \textbf{master}, converts the pulse code modulated (PCM) data from memory via DMA into pulse density modulation (PDM) data, and then sends the data out. Used signals: WS and DATA. PDM standard is supported in this mode by user configuration.\\
\end{longtable}

\section{Features}

The I2S module has the following features:

\begin{itemize}
    \item Master mode and slave mode
    \item Full-duplex and half-duplex communications
    \item Separate TX and RX units that can work independently or simultaneously
    \item A variety of audio standards supported: 
        \begin{itemize}
            \item TDM Philips standard
            \item TDM MSB alignment standard
            \item TDM PCM standard
            \item PDM standard
        \end{itemize}
    \item Configurable high-precision sample clock:
        \begin{itemize}
            \item Supports the following frequencies: 8 kHz, 16 kHz, 32 kHz, 44.1 kHz, 48 kHz, 88.2 kHz, 96 kHz, 128 kHz, and 192 kHz (Note that in slave mode, due to the frequency limitation of the clock source, the maximum sampling frequency is limited by the data bit width and number of channels. For detailed information, refer to Section \ref{sec:i2s-tx-rx-clk})
        \end{itemize}
    \item 8-/16-/24-/32-bit data communication
    \item Direct Memory Access (DMA)
    \item Standard I2S interface interrupts 
\end{itemize}

\section{System Architecture}
\begin{figure}[H]
    \centering
    \fbox{\includegraphics[width=0.9\textwidth]{\modulefiles/figures/i2s_graph_20221121}}
    \caption{\chipname{} I2S System Diagram}
    \label{fig:i2s-sys-diagram}
\end{figure}

Figure \ref{fig:i2s-sys-diagram} shows the structure of \chipname{} I2S module, consisting of:
\begin{itemize}
    \item Transmit control unit (TX Unit)
    \item Receive control unit (RX Unit)
    \item Input and output timing unit (I/O Sync)
    \item Clock divider (Clock Generator)
    \item 64 x 32-bit TX FIFO
    \item 64 x 32-bit RX FIFO
    \item Compress/Decompress units
\end{itemize}
I2S module supports direct memory access (DMA) to internal memory. For more information, see Chapter \ref{mod:dma} \textit{\nameref{mod:dma}}. 

Both the TX unit and the RX unit have a three-line interface that uses a bit clock line (BCK), a word select line (WS), and a serial data line (SD). The SD line of the TX unit is fixed as output, and the SD line of the RX unit as input. BCK and WS signal lines for TX unit and RX unit can be configured as master output mode or slave input mode.

The signal bus of I2S module is shown at the right part of Figure \ref{fig:i2s-sys-diagram}. The naming of these signals in RX and TX units follows the pattern: I2S\textcolor{red}{A}\_\textcolor{red}{B}\_\textcolor{red}{C}, such as I2S\textcolor{red}{I}\_\textcolor{red}{BCK}\_\textcolor{red}{in}.

\begin{itemize}
    \item ``A” represents the direction of data bus, which includes:
    \begin{itemize}
        \item ``I”: input, receiving
        \item ``O”: output, transmitting
    \end{itemize}
    \item ``B” represents the signal function, which includes:
    \begin{itemize}
        \item BCK
        \item WS
        \item SD
    \end{itemize}
    \item ``C” represents the signal direction, which includes:
    \begin{itemize}
        \item “in”: input signal into I2S module
        \item ``out”: output signal from I2S module
    \end{itemize}
\end{itemize}

Table \ref{tab:i2s-signal-descr} provides a detailed description of I2S signals.


\begin{table}[H]
    \centering
    \caption{I2S Signal Description}
    \label{tab:i2s-signal-descr}
    \begin{tabular}{|p{2.8cm}|p{1.8cm}|p{10.2cm}| }
    \hline
    \rowcolor{lightgray}
    \textbf{Signal} & \textbf{Direction} & \textbf{Function} \\
    \hline
    I2SI\_BCK\_in  & Input & In I2S slave mode, inputs BCK signal for RX unit. \\
    \hline
    I2SI\_BCK\_out& Output & In I2S master mode, outputs BCK signal for RX unit. \\
    \hline
    I2SI\_WS\_in & Input & In I2S slave mode, inputs WS signal for RX unit. \\
    \hline
    I2SI\_WS\_out& Output & In I2S master mode, outputs WS signal for RX unit. \\\hline
    I2SI\_Data\_in& Input & Works as the serial input data bus for I2S RX unit.\\\hline
    I2SO\_Data\_out& Output & Works as the serial output data bus for I2S TX unit.\\\hline
    
    I2SO\_BCK\_in&  Input & In I2S slave mode, inputs BCK signal for TX unit. \\\hline
    I2SO\_BCK\_out&  Output & In I2S master mode, outputs BCK signal for TX unit. \\\hline
    I2SO\_WS\_in&  Input & In I2S slave mode, inputs WS signal for TX unit. \\\hline
    I2SO\_WS\_out& Output & In I2S master mode, outputs WS signal for TX unit. \\\hline
    I2S\_MCLK\_in& Input & In I2S slave mode, works as a clock source from the external master.\\\hline
    I2S\_MCLK\_out& Output & In I2S master mode, works as a clock source for the external slave. \\\hline
    \end{tabular}
\end{table}

\vspace{-0.5em}
\begin{tiplisting}
Any required signals of I2S must be mapped to the chip's pins via GPIO matrix. For more information, see Chapter \ref{mod:iomuxgpio} \textit{\nameref{mod:iomuxgpio}}.
\end{tiplisting}

\section{Supported Audio Standards} \label{sec:i2s-supported-audio-standards}


\chipname{} I2S supports multiple audio standards, including TDM Philips standard, TDM MSB alignment standard, TDM PCM standard, and PDM standard.

Select the needed standard by configuring the following bits:
\begin{itemize}
    \item \hyperref[fielddesc:I2SRXMSBSHIFT]{I2S\_TX/RX\_TDM\_EN}
    \begin{itemize}
        \item 0: disable TDM mode.
        \item 1: enable TDM mode.
    \end{itemize}
    \item \hyperref[fielddesc:I2SRXMSBSHIFT]{I2S\_TX/RX\_PDM\_EN}
    \begin{itemize}
        \item 0: disable PDM mode.
        \item 1: enable PDM mode.
    \end{itemize}    
    \item \hyperref[fielddesc:I2SRXMSBSHIFT]{I2S\_TX/RX\_MSB\_SHIFT}
    \begin{itemize}
        \item 0: WS and SD signals change simultaneously, i.e., enable MSB alignment standard.  
        \item 1: WS signal changes one BCK clock cycle earlier than SD signal, i.e., enable Philips standard or select PCM standard.
    \end{itemize}
    \item \hyperref[fielddesc:I2SRXPCMBYPASS]{I2S\_TX/RX\_PCM\_BYPASS}
    \begin{itemize}
        \item 0: enable PCM standard.
        \item 1: disable PCM standard.
    \end{itemize}
\end{itemize}

\subsection{TDM Philips Standard}


Philips specifications require that WS signal changes one BCK clock cycle earlier than SD signal on BCK falling edge, which means that WS signal is valid from one clock cycle before transmitting the first bit of channel data and changes one clock before the end of channel data transfer. SD signal line transmits the most significant bit of audio data first. 

Compared with Philips standard, TDM Philips standard supports multiple channels. See Figure \ref{fig:i2s-tdm-philips-std-timing-diagram}.

\begin{figure}[H]
    \centering
    \includegraphics[width=1.0\textwidth]{\modulefiles/figures/i2s_tdm_phillips_mode.png}
    \caption{TDM Philips Standard Timing Diagram}
    \label{fig:i2s-tdm-philips-std-timing-diagram}
\end{figure}




\subsection{TDM MSB Alignment Standard}

MSB alignment specifications require that WS and SD signals change simultaneously on the falling edge of BCK. The WS signal is valid until the end of channel data transfer. The SD signal line transmits the most significant bit of audio data first.

Compared with MSB alignment standard, TDM MSB alignment standard supports multiple channels. See Figure \ref{fig:i2s-tdm-msb-align-std-timing-diagram}.

\begin{figure}[H]
    \centering
    \includegraphics[width=1.0\textwidth]{\modulefiles/figures/i2s_tdm_msb_mode.png}
    \caption{TDM MSB Alignment Standard Timing Diagram}
    \label{fig:i2s-tdm-msb-align-std-timing-diagram}
\end{figure}



\subsection{TDM PCM Standard}

Short frame synchronization under PCM standard requires that WS signal changes one BCK clock cycle earlier than SD signal on the falling edge of BCK, which means that the WS signal becomes valid one clock cycle before transferring the first bit of channel data and remains unchanged in this BCK clock cycle. SD signal line transmits the most significant bit of audio data first. 

Compared with PCM standard, TDM PCM standard supports multiple channels. See Figure \ref{fig:i2s-tdm-pcm-std-timing-diagram}.

\begin{figure}[H]
    \centering
    \includegraphics[width=1.0\textwidth]{\modulefiles/figures/i2s_tdm_pcm_mode.png}
    \caption{TDM PCM Standard Timing Diagram}
    \label{fig:i2s-tdm-pcm-std-timing-diagram}
\end{figure}



\subsection{PDM Standard}
Under PDM standard, WS signal changes continuously during data transmission. The low-level and high-level of this signal indicates the left channel and right channel respectively. WS and SD signals change simultaneously on the falling edge of BCK. See Figure \ref{fig:i2s-pdm-std-timing-diagram}. 

\begin{figure}[H]
    \centering
    \includegraphics[width=1.0\textwidth]{\modulefiles/figures/i2s_pdm_mode.png}
    \caption{PDM Standard Timing Diagram}
    \label{fig:i2s-pdm-std-timing-diagram}
\end{figure}


\section{I2S TX/RX Clock}\label{sec:i2s-tx-rx-clk}

I2S\_TX/RX\_CLK is the master clock of I2S TX/RX unit, divided from:
\begin{itemize}
    \item 40 MHz XTAL\_CLK 
    \item 160 MHz PLL\_F160M\_CLK
    \item 240 MHz PLL\_F240M\_CLK
    \item or external input clock: I2S\_MCLK\_in 
\end{itemize}

The serial clock (BCK) of the I2S TX/RX unit is divided from I2S\_TX/RX\_CLK, as shown in Figure \ref{fig:i2s-clk}. \hyperref[fielddesc:PCRI2SRXCLKMSEL]{PCR\_I2S\_TX/RX\_CLKM\_SEL} is used to select clock source for TX/RX unit, and \hyperref[fielddesc:PCRI2SRXCLKMEN]{PCR\_I2S\_TX/RX\_CLKM\_EN} to enable or disable the clock source.

\begin{figure}[H]
    \centering
    \fbox{\includegraphics[width=0.6\textwidth]{\modulefiles/figures/i2s_clk_20221122}}
    \caption{I2S Clock Generator}
    \label{fig:i2s-clk}
\end{figure}


The following formula shows the relation between I2S\_TX/RX\_CLK frequency $f_{\textrm{I2S\_TX/RX\_CLK}}$ and the divider clock source frequency $f_{\textrm{I2S\_CLK\_S}}$:
$$f_{\textrm{I2S\_TX/RX\_CLK}}=\frac{f_{\textrm{I2S\_CLK\_S}}}{\textrm{N} + \frac{\textrm{b}}{\textrm{a}}}$$

N is an integer value between 2 and 256. The value of N corresponds to the value of  \hyperref[fielddesc:PCRI2STXCLKMDIVNUM]{PCR\_I2S\_TX/RX\_CLKM\_DIV\_NUM} in register  \hyperref[regdesc:PCRI2STXCLKMCONFREG]{PCR\_I2S\_TX/RX\_CLKM\_CONF\_REG} as follows:

\begin{itemize}
    \item When  \hyperref[fielddesc:PCRI2STXCLKMDIVNUM]{PCR\_I2S\_TX/RX\_CLKM\_DIV\_NUM} = 0, N = 256;
    \item When  \hyperref[fielddesc:PCRI2STXCLKMDIVNUM]{PCR\_I2S\_TX/RX\_CLKM\_DIV\_NUM} = 1, N = 2;
    \item When  \hyperref[fielddesc:PCRI2STXCLKMDIVNUM]{PCR\_I2S\_TX/RX\_CLKM\_DIV\_NUM} has any other value, N =  \hyperref[fielddesc:PCRI2STXCLKMDIVNUM]{PCR\_I2S\_TX/RX\_CLKM\_DIV\_NUM}.

\end{itemize}

The values of ``a" and ``b" in fractional divider depend only on x, y, z, and yn1. The corresponding formulas are as follows:

\begin{itemize}
    \item When $\textrm{b} <= \frac{\textrm{a}}{2}$, $\textrm{yn1} = 0$, $\textrm{x} = floor([\frac{\textrm{a}}{\textrm{b}}]) - 1$, $\textrm{y} = \textrm{a}\%\textrm{b}$, $\textrm{z} = \textrm{b}$; 
    \item When $\textrm{b} > \frac{\textrm{a}}{2}$, $\textrm{yn1}=1$, $\textrm{x} = floor([\frac{\textrm{a}}{\textrm{a - b}}]) - 1$, $\textrm{y} = \textrm{a}\%\textrm{(a - b)}$, $\textrm{z}=\textrm{a - b}$.
\end{itemize}

The values of x, y, z, and yn1 are configured in \hyperref[fielddesc:PCRI2STXCLKMDIVX]{PCR\_I2S\_TX/RX\_CLKM\_DIV\_X}, \hyperref[fielddesc:PCRI2STXCLKMDIVY]{PCR\_I2S\_TX/RX\_CLKM\_DIV\_Y}, \hyperref[fielddesc:PCRI2STXCLKMDIVZ]{PCR\_I2S\_TX/RX\_CLKM\_DIV\_Z}, and \hyperref[fielddesc:PCRI2STXCLKMDIVYN1]{PCR\_I2S\_TX/RXCLKM\_DIV\_YN1}.

To configure the integer divider, clear \hyperref[fielddesc:PCRI2STXCLKMDIVX]{PCR\_I2S\_TX/RX\_CLKM\_DIV\_X} and \hyperref[fielddesc:PCRI2STXCLKMDIVZ]{PCR\_I2S\_TX/RX\_CLKM\_DIV\_Z}, then set \hyperref[fielddesc:PCRI2STXCLKMDIVY]{PCR\_I2S\_TX/RX\_CLKM\_DIV\_Y} to 1.


\begin{tiplisting}
Using fractional divider may introduce some clock jitter.
\end{tiplisting}


In master TX mode, the serial clock BCK for I2S TX unit is I2SO\_BCK\_out divided from I2S\_TX\_CLK, which is:

$$f_{\textrm{I2SO\_BCK\_out}} = \frac{f_{\textrm{I2S\_TX\_CLK}}}{\textrm{MO}}$$


``MO” is an integer value:
$$ \textrm{MO} = \textrm{\hyperref[fielddesc:I2STXBCKDIVNUM]{I2S\_TX\_BCK\_DIV\_NUM}} + 1 $$

\vspace{-1em}
\begin{tiplisting}
Note that \hyperref[fielddesc:I2STXBCKDIVNUM]{I2S\_TX\_BCK\_DIV\_NUM} must not be configured as 1.
\end{tiplisting}

In master RX mode, the serial clock BCK for I2S RX unit is I2SI\_BCK\_out divided from I2S\_RX\_CLK, which is:

$$f_{\textrm{I2SI\_BCK\_out}} = \frac{f_{\textrm{I2S\_RX\_CLK}}}{\textrm{MI}}$$

``MI” is an integer value:
$$ \textrm{MI} = \textrm{\hyperref[fielddesc:I2STXBCKDIVNUM]{I2S\_RX\_BCK\_DIV\_NUM}} + 1 $$

\clearpage

\begin{tiplisting}
\vspace{-2em}
\begin{itemize}
    \item \hyperref[fielddesc:I2SRXBCKDIVNUM]{I2S\_RX\_BCK\_DIV\_NUM} must not be configured as 1.
    \item In I2S slave mode, make sure $f_{\textrm{I2S\_TX/RX\_CLK}}$ >= 8 * $f_{\textrm{BCK}}$. The I2S module can output I2S\_MCLK\_out as the master clock for peripherals.
\end{itemize}
\end{tiplisting}

\section{I2S Reset} \label{sec:i2s-reset}

The units and FIFOs in I2S module are reset by the following bits.
\begin{itemize}
\item I2S TX/RX units: reset by the bits  \hyperref[fielddesc:I2STXRESET]{I2S\_TX\_RESET} and  \hyperref[fielddesc:I2SRXRESET]{I2S\_RX\_RESET};
\item I2S TX/RX FIFO: reset by the bits  \hyperref[fielddesc:I2STXFIFORESET]{I2S\_TX\_FIFO\_RESET} and  \hyperref[fielddesc:I2SRXFIFORESET]{I2S\_RX\_FIFO\_RESET}.
\end{itemize}

\begin{tiplisting}
The I2S module clock must be configured first before the module and FIFO are reset.
\end{tiplisting} 

\section{I2S Master/Slave Mode}

The \chipname{} I2S module can operate as a master or a slave in half-duplex and full-duplex communications, depending on the configuration of  \hyperref[fielddesc:I2SRXSLAVEMOD]{I2S\_RX\_SLAVE\_MOD} and  \hyperref[fielddesc:I2STXSLAVEMOD]{I2S\_TX\_SLAVE\_MOD}.

\begin{itemize}
    \item \hyperref[fielddesc:I2STXSLAVEMOD]{I2S\_TX\_SLAVE\_MOD} 
    \begin{itemize}
        \item 0: master TX mode
        \item 1: slave TX mode
    \end{itemize}
    \item \hyperref[fielddesc:I2SRXSLAVEMOD]{I2S\_RX\_SLAVE\_MOD}
    \begin{itemize}
        \item 0: master RX mode
        \item 1: slave RX mode
    \end{itemize}
\end{itemize}

\subsection{Master/Slave TX Mode} \label{sec:i2s-master-slave-tx-mode}
\begin{itemize}
    \item I2S works as a master transmitter:
    \begin{itemize}
        \item Set  \hyperref[fielddesc:I2STXSTART]{I2S\_TX\_START} to start transmitting data.
        \item TX unit keeps driving the clock signal and serial data.
        \item If  \hyperref[fielddesc:I2STXSTOPEN]{I2S\_TX\_STOP\_EN} is set and all the data in FIFO is transmitted, the master stops transmitting data and clock signals.
        \item If  \hyperref[fielddesc:I2STXSTOPEN]{I2S\_TX\_STOP\_EN} is cleared and all the data in FIFO is transmitted, meanwhile no new data is filled into FIFO, then the TX unit keeps sending the last data frame and clock signal.
        \item Master stops sending data when the bit  \hyperref[fielddesc:I2STXSTART]{I2S\_TX\_START} is cleared.
    \end{itemize}
    \item I2S works as a slave transmitter:
     \begin{itemize}
        \item Set  \hyperref[fielddesc:I2STXSTART]{I2S\_TX\_START}.
        \item Wait for the master BCK clock to enable a transmit operation.
        \item If  \hyperref[fielddesc:I2STXSTOPEN]{I2S\_TX\_STOP\_EN} is set and all the data in FIFO is transmitted, then the slave keeps sending zeros, till the master stops providing BCK signal.
        \item If  \hyperref[fielddesc:I2STXSTOPEN]{I2S\_TX\_STOP\_EN} is cleared and all the data in FIFO is transmitted, meanwhile no new data is filled into FIFO, then the TX unit keeps sending the last data frame. 
        \item If \hyperref[fielddesc:I2STXSTART]{I2S\_TX\_START} is cleared, slave keeps sending zeros till the master stops providing BCK clock signal.
    \end{itemize}
\end{itemize}

\subsection{Master/Slave RX Mode}
\begin{itemize}
    \item I2S works as a master receiver:
    \begin{itemize}
        \item Set  \hyperref[fielddesc:I2SRXSTART]{I2S\_RX\_START} to start receiving data.
        \item RX unit keeps outputting clock signal and sampling input data.
        \item RX unit stops receiving data when the bit  \hyperref[fielddesc:I2SRXSTART]{I2S\_RX\_START} is cleared.
    \end{itemize}
    \item I2S works as a slave receiver:
     \begin{itemize}
        \item Set  \hyperref[fielddesc:I2SRXSTART]{I2S\_RX\_START}.
        \item Wait for master BCK signal to start receiving data.
        \item RX unit stops receiving data when the bit \hyperref[fielddesc:I2SRXSTART]{I2S\_RX\_START} is cleared.
    \end{itemize}
\end{itemize}

\section{Transmitting Data}\label{sec:i2s-tx-data}

\vspace{1em}
\begin{tiplisting}
Updating the configuration described in this and subsequent sections requires to set \hyperref[fielddesc:I2STXUPDATE]{I2S\_TX\_UPDATE} accordingly to synchronize registers from APB clock domain to TX clock domain. For more detailed configuration, see Section \ref{sec:i2s-configure-i2s-tx-mode}.
\end{tiplisting}

\vspace{-1em}
In TX mode, I2S first reads data through DMA and sends these data out via output signals according to the configured data mode and channel mode.
\subsection{Data Format Control}\label{sec:i2s-data-control}
Data format is controlled in the following phases:
\begin{itemize}
    \item Phase I: read data from memory and write it to TX FIFO;
    \item Phase II: read the data to send (TX data) from TX FIFO and convert the data according to the output data mode;
    \item Phase III: clock out the TX data serially.
\end{itemize}

\subsubsection{Bit Width Control of Channel Valid Data}
The bit width of valid data in each channel is determined by \hyperref[fielddesc:I2STXBITSMOD]{I2S\_TX\_BITS\_MOD} and  \hyperref[fielddesc:I2STX24FILLEN]{I2S\_TX\_24\_FILL\_EN}. For details, see the table below.

\begin{table}[H]
    \centering
    \caption{Bit Width of Channel Valid Data}
    \label{tab:i2s-bit-width-chan-valid-data}
    \begin{threeparttable}
    \begin{tabular}{|p{4.2cm}|R{3.1cm}|R{3.5cm}|}
    \hline
    \rowcolor{lightgray}
    \textbf{Channel Valid Data Width} & \textbf{\hyperref[fielddesc:I2STXBITSMOD]{I2S\_TX\_BITS\_MOD}}  & \textbf{\hyperref[fielddesc:I2STX24FILLEN]{I2S\_TX\_24\_FILL\_EN}} \\ \hline
    \multirow{2}*{32} & 31 & x\tnote{*} \\\cline{2-3}
                      & 23 & 1 \\\hline

                24    & 23 & 0 \\\hline

                16     & 15 & x \\\hline
                8     & 7 & x \\\hline
    \end{tabular}
            \begin{tablenotes}
            \item[*] ``x" represents that this value is ignored.
        \end{tablenotes}

    \end{threeparttable}
\end{table}

\subsubsection{Endian Control of Channel Valid Data}
When I2S reads data through DMA, the data endian under various data width is controlled by \hyperref[fielddesc:I2STXBIGENDIAN]{I2S\_TX\_BIG\_ENDIAN}. Table \ref{tab:i2s-endian-chan-valid-data} shows how \hyperref[fielddesc:I2STXBIGENDIAN]{I2S\_TX\_BIG\_ENDIAN} controls the data reading with different channel valid data widths.
\begin{table}[H]
    \centering
    \caption{Endian of Channel Valid Data}
    \label{tab:i2s-endian-chan-valid-data}
    \begin{tabular}{|p{4.2cm}|p{2.7cm}|p{4.2cm}|R{3.5cm}|}
    \hline
    \rowcolor{lightgray}
    \textbf{Channel Valid Data Width} & \textbf{Original Data}  & \textbf{Endian of Processed Data}  & \textbf{\hyperref[fielddesc:I2STXBIGENDIAN]{I2S\_TX\_BIG\_ENDIAN}} \\ \hline
    \multirow{2}*{32} & \multirow{2}*{\{B3, B2, B1, B0\}} & \{B3, B2, B1, B0\} & 0 \\\cline{3-4}
                      && \{B0, B1, B2, B3\} & 1 \\\hline

    \multirow{2}*{24} & \multirow{2}*{\{B2, B1, B0\}} & \{B2, B1, B0\} & 0 \\\cline{3-4}
                      && \{B0, B1, B2\} & 1 \\\hline

    \multirow{2}*{16} & \multirow{2}*{\{B1, B0\}} & \{B1, B0\} & 0 \\\cline{3-4}
                      && \{B0, B1\} & 1 \\\hline

                8     & \{B0\} & \{B0\} & x \\\hline
    \end{tabular}
\end{table}

\begin{tiplisting}
B0, B1, B2, B3 each represents an 8-bit data, and the symbol \{\} means that the bytes are combined together. For example, \{B3, B2, B1, B0\} represents a 32-bit number, wherein B0 represents bit 0-7, B1 represents bit 8-15, B2 represents bit 16-23, and B3 represents bit 24-31.
\end{tiplisting}


\subsubsection{A-law/$\mu$-law Compression and Decompression}
\chipname{} I2S compresses/decompresses the valid data into 32-bit by A-law or by $\mu$-law. If the bit width of valid data is smaller than 32, zeros are filled to the extra high bits of the data to be compressed/decompressed by default.
\begin{tiplisting}
Extra high bits here mean the bits[31: channel valid data width] of the data to be compressed/decompressed.
\end{tiplisting}

Configure \hyperref[fielddesc:I2STXPCMBYPASS]{I2S\_TX\_PCM\_BYPASS}:
\begin{itemize}
\item 0: compress or decompress the data
\item 1: do not compress or decompress the data
\end{itemize}
Configure \hyperref[fielddesc:I2STXPCMCONF]{I2S\_TX\_PCM\_CONF}: 
\begin{itemize}
\item 0: decompress the data using A-law
\item 1: compress the data using A-law
\item 2: decompress the data using $\mu$-law
\item 3: compress the data using $\mu$-law
\end{itemize}
At this point, the first phase of data format control is completed.

\subsubsection{Bit Width Control of Channel TX Data} 
The TX data width in each channel is determined by \hyperref[fielddesc:I2STXTDMCHANBITS]{I2S\_TX\_TDM\_CHAN\_BITS}.

\begin{itemize}
\item If TX data width in each channel is larger than the valid data width, zeros will be filled to these extra bits. Configure  \hyperref[fielddesc:I2STXLEFTALIGN]{I2S\_TX\_LEFT\_ALIGN}:
\begin{itemize}
\item 0: the valid data is at the lower bits of TX data. Zeros are filled into higher bits of TX data;
\item 1: the valid data is at the higher bits of TX data. Zeros are filled into lower bits of TX data.
\end{itemize}
\item If the TX data width in each channel is smaller than the valid data width, only the lower bits of valid data are sent out, and the higher bits are discarded.
\end{itemize}
At this point, the second phase of data format control is completed.

\subsubsection{Bit Order Control of Channel Data}
The data bit order in each channel is controlled by \hyperref[fielddesc:I2STXBITORDER]{I2S\_TX\_BIT\_ORDER}:
\begin{itemize}
\item 0: Not reverse the valid data bit order;
\item 1: Reverse the valid data bit order.
\end{itemize}
At this point, the data format control is completed. The data after format control will be sent sequentially from high to low. Figure \ref{fig:i2s-tx-data-format-control} shows the complete process of TX data format control.


\begin{figure}[H]
    \centering
    \includegraphics[width=0.8\textwidth]{\modulefiles/figures/i2s_tx_data_config.png}
    \caption{TX Data Format Control}
    \label{fig:i2s-tx-data-format-control}
\end{figure}

\subsection{Channel Mode Control}
\chipname{} I2S supports both TDM TX mode and PDM TX mode. Set \hyperref[fielddesc:I2STXTDMEN]{I2S\_TX\_TDM\_EN} to enable TDM TX mode, or set \hyperref[fielddesc:I2STXPDMEN]{I2S\_TX\_PDM\_EN} to enable PDM TX mode.

\begin{tiplisting}
\vspace{-2em}
\begin{itemize}
    \item \hyperref[fielddesc:I2STXTDMEN]{I2S\_TX\_TDM\_EN} and \hyperref[fielddesc:I2STXPDMEN]{I2S\_TX\_PDM\_EN} must not be cleared or set simultaneously. 
    \item Most stereo I2S codecs can be controlled by setting the I2S module into 2-channel mode under TDM standard.
    \end{itemize}
\end{tiplisting}

\subsubsection{I2S Channel Control in TDM TX Mode}\label{sec:i2s-chan-control-tdm-mode} 
In TDM TX mode, the total number of TX channels supported is related to the channel valid data width for I2S as follows:

\begin{table}[H]
    \centering
    \caption{The Matching Between Valid Data Width and Number of TX Channel Supported}
    \label{tab:i2s-tx-support-channel}
    \begin{tabular}{|l|r|}
    \hline
    \rowcolor{lightgray}
    \textbf{Channel Valid Data Width} & \textbf{Total Number of Channels Supported} \\ \hline
        32  & 4     \\\hline
        24  & 5     \\\hline
        16  & 8     \\\hline
        8   & 16    \\\hline
    \end{tabular}
\end{table}

The total number of TX channels in use is controlled by \hyperref[fielddesc:I2STXTDMTOTCHANNUM]{I2S\_TX\_TDM\_TOT\_CHAN\_NUM}. For example, if  \hyperref[fielddesc:I2STXTDMTOTCHANNUM]{I2S\_TX\_TDM\_TOT\_CHAN\_NUM} is set to 5, six channels in total (channel 0 \verb+~+ 5) will be used to transmit data. See Figure \ref{fig:i2s-tdm-chan-control}.

In these TX channels, if \hyperref[fielddesc:I2STXTDMCHAN0EN]{I2S\_TX\_TDM\_CHAN{\textcolor{red}{n}}\_EN} is set to:

\begin{itemize}
    \item 1: this channel sends the channel data out;
    \item 0: the TX data to be sent by this channel is controlled by \hyperref[fielddesc:I2STXCHANEQUAL]{I2S\_TX\_CHAN\_EQUAL}:
    \begin{itemize}
        \item 1: the data of previous channel is sent out;
        \item 0: the data stored in \hyperref[fielddesc:I2SSINGLEDATA]{I2S\_SINGLE\_DATA} is sent out.
    \end{itemize}
\end{itemize}
In TDM TX  master mode, WS signal is controlled by \hyperref[fielddesc:I2STXWSIDLEPOL]{I2S\_TX\_WS\_IDLE\_POL} and \hyperref[fielddesc:I2STXTDMWSWIDTH]{I2S\_TX\_TDM\_WS\_WIDTH}:
\begin{itemize}
    \item \hyperref[fielddesc:I2STXWSIDLEPOL]{I2S\_TX\_WS\_IDLE\_POL}: the default level of WS signal;
    \item \hyperref[fielddesc:I2STXTDMWSWIDTH]{I2S\_TX\_TDM\_WS\_WIDTH}: the cycles the WS default level lasts for when transmitting all channel data.
\end{itemize}
\hyperref[fielddesc:I2STXHALFSAMPLEBITS]{I2S\_TX\_HALF\_SAMPLE\_BITS} x 2 is equal to the BCK cycles in one WS period.

\textbf{TDM Channel Configuration Example}

In this example, register configuration is as follows:
\begin{itemize}
    \item I2S\_TX\_TDM\_CHAN\_NUM = 5, i.e., channel 0 \verb+~+ 5 are used to transmit data.
    \item I2S\_TX\_CHAN\_EQUAL = 1, i.e., that data of previous channel will be transmitted if the bit I2S\_TX\_TDM\_CHAN\regindex{n}\_EN is cleared. \regindex{n} = 0 \verb+~+ 5.
    \item I2S\_TX\_TDM\_CHAN0/2/5\_EN = 1, i.e., these channels send their channel data out.
    \item I2S\_TX\_TDM\_CHAN1/3/4\_EN = 0, i.e., these channels send the previous channel data out.
\end{itemize}

Once the configuration is done, data is transmitted as follows.

\begin{figure}[H]
    \centering
    \includegraphics[width=1.0\textwidth]{\modulefiles/figures/i2s_tdm_channel_config.pdf}
    \caption{TDM Channel Control}
    \label{fig:i2s-tdm-chan-control}
\end{figure}

\subsubsection{I2S Channel Control in PDM TX Mode} 

\chipname{} I2S supports two PDM TX modes, namely, normal PDM TX mode and PCM-to-PDM TX mode.

In PDM TX mode, fetching data through DMA is controlled by \hyperref[fielddesc:I2STXMONO]{I2S\_TX\_MONO} and \hyperref[fielddesc:I2STXMONOFSTVLD]{I2S\_TX\_MONO\_FST\_VLD}. See Table \ref{tab:i2s-data-fetch-control-pdm-mode}. Please configure the two bits according to the data stored in memory, be it the single-channel or dual-channel data. 
\begin{table}[H]
    \centering
    \caption{Data-Fetching Control in PDM Mode}
    \label{tab:i2s-data-fetch-control-pdm-mode}
    \begin{tabular}{|p{6.6cm}|l|r|r|}
    \hline
    \rowcolor{lightgray}
    \textbf{Data-Fetching Control Option} &\textbf{Mode} &\textbf{\hyperref[fielddesc:I2STXMONO]{I2S\_TX\_MONO}}  & \textbf{\hyperref[fielddesc:I2STXMONOFSTVLD]{I2S\_TX\_MONO\_FST\_VLD}} \\ \hline
                        Post data-fetching request to DMA at any edge of WS signal & Stereo mode & 0 & x        \\ \hline
                        Post data-fetching request to DMA only at the second half period of WS signal  & Mono mode & 1 & 0   \\ \hline
                        Post data-fetching request to DMA only at the first half period of WS signal & Mono mode & 1 & 1      \\ \hline
    \end{tabular}
\end{table}

When the I2S is in PDM TX master mode, the default level of WS signal is controlled by \hyperref[fielddesc:I2STXWSIDLEPOL]{I2S\_TX\_WS\_IDLE\_POL}, and the WS signal frequency is half of the BCK signal frequency. The configuration of WS signal is similar to that of BCK signal. Please refer to Section \ref{sec:i2s-tx-rx-clk} and Figure \ref{sec:i2s-tx-rx-clk}.

In \textbf{normal PDM TX mode}, I2S channel mode is controlled by  \hyperref[fielddesc:I2STXCHANMOD]{I2S\_TX\_CHAN\_MOD} and  \hyperref[fielddesc:I2STXWSIDLEPOL]{I2S\_TX\_WS\_IDLE\_POL}. See the table below. 

\begin{table}[H]
    \centering
    \caption{I2S Channel Control in Normal PDM TX Mode}
    \label{tab:i2s-chan-control-pdm-mode}
    \begin{threeparttable}
    \begin{tabular}{|p{2.4cm}|L{4.7cm}|L{4.7cm}|R{1.3cm}|R{1.4cm}|}
    \hline
    \rowcolor{lightgray}
    \textbf{Channel Control Option} & \textbf{Left Channel} &\textbf{Right Channel} & \textbf{Mode Control Field}\tnote{1} & \textbf{Channel Select Bit}\tnote{2}\\ \hline
    Stereo mode & Transmit the left channel data & Transmit the right channel data &0 & x \\\hline
                         \multirow{8}*{Mono mode}  & Transmit the left channel data & Transmit the left channel data & 1 & 0         \\ \cline{2-5}
                                           & Transmit the right channel data & Transmit the right channel data& 1 & 1\\ \cline{2-5}
                                           & Transmit the right channel data & Transmit the right channel data & 2 & 0 \\ \cline{2-5}
                                           & Transmit the left channel data & Transmit the left channel data &2 & 1 \\ \cline{2-5}
                                           & Transmit the value of ``single”\tnote{3} & Transmit the right channel data & 3 & 0 \\ \cline{2-5}
                                           & Transmit the left channel data & Transmit the value of ``single” &3 & 1 \\ \cline{2-5}
                                           & Transmit the left channel data & Transmit the value of ``single”&4 & 0 \\ \cline{2-5}
                                           & Transmit the value of ``single” & Transmit the right channel data & 4 & 1 \\ \hline
    \end{tabular}
            \begin{tablenotes}
            \item[1] \hyperref[fielddesc:I2STXCHANMOD]{I2S\_TX\_CHAN\_MOD}
            \item[2] \hyperref[fielddesc:I2STXWSIDLEPOL]{I2S\_TX\_WS\_IDLE\_POL}
            \item[3] The ``single” value is equal to the value of  \hyperref[fielddesc:I2SSINGLEDATA]{I2S\_SINGLE\_DATA}.
        \end{tablenotes}

    \end{threeparttable}
\end{table}

In \textbf{PCM-to-PDM TX mode}, the PCM data through DMA is converted to PDM data and then output in PDM signal format. Configure  \hyperref[fielddesc:I2SPCM2PDMCONVEN]{I2S\_PCM2PDM\_CONV\_EN} to enable this mode. The register configuration for PCM-to-PDM TX mode is as follows:

\begin{itemize}
    \item Configure 1-line PDM output format or 1-/2-line DAC output mode as the table below:

\begin{table}[H]
    \centering
    \caption{PCM-to-PDM TX Mode}
    \label{tab:i2s-pcm-pdm-output-mode}
    \begin{tabular}{|p{4.3cm}|R{5.0cm}|R{4.8cm}|}
    \hline
    \rowcolor{lightgray}
    \textbf{Channel Output Format}&\textbf{\hyperref[fielddesc:I2STXPDMDACMODEEN]{I2S\_TX\_PDM\_DAC\_MODE\_EN}} & \textbf{\hyperref[fielddesc:I2STXPDMDAC2OUTEN]{I2S\_TX\_PDM\_DAC\_2OUT\_EN}}\\ \hline
    1-line PDM output format\hyperref[note-1]{$^1$}& 0 & x \\ \hline
    1-line DAC output format\hyperref[note-2]{$^2$} &1   & 0  \\ \hline 
    2-line DAC output format & 1 & 1 \\ \hline
    \end{tabular}
\end{table}

\begin{tiplisting}
\vspace{-2em}
\begin{enumerate}
    \item \label{note-1} In PDM output format, SD data of two channels is sent out in one WS period.
    \item \label{note-2} In DAC output format, SD data of one channel is sent out in one WS period.
\end{enumerate}
\end{tiplisting}
    
    \item Configure sampling frequency and upsampling rate: \newline
    In PCM-to-PDM TX mode, PDM clock frequency is equal to BCK frequency. The relation of sampling frequency ($f_{\textrm{Sampling}}$) and BCK frequency is as follows:
    $$f_{\textrm{Sampling}} = \frac{f_{\textrm{BCK}}}{\textrm{OSR}}$$
    
    Upsampling rate (OSR) is related to \hyperref[fielddesc:I2STXPDMSINCOSR2]{I2S\_TX\_PDM\_SINC\_OSR2} as follows:
    $${\textrm{OSR}} = {\textrm{\hyperref[fielddesc:I2STXPDMSINCOSR2]{I2S\_TX\_PDM\_SINC\_OSR2}}} \times {\textrm{64}}$$
    
    Sampling frequency $f_{\textrm{Sampling}}$ is related to \hyperref[fielddesc:I2STXPDMFS]{I2S\_TX\_PDM\_FS} as follows:
    $$f_{\textrm{Sampling}} = {\textrm{\hyperref[fielddesc:I2STXPDMFS]{I2S\_TX\_PDM\_FS}}} \times {\textrm{100}}$$
    
    Configure the registers according to needed sampling frequency, upsampling rate, and PDM clock frequency.
\end{itemize}


\textbf{PDM Channel Configuration Example}

In this example, the register configuration is as follows.
\begin{itemize}
    \item \hyperref[fielddesc:I2SPCM2PDMCONVEN]{I2S\_PCM2PDM\_CONV\_EN} = 0, i.e., the normal PDM TX mode is selected.
    \item \hyperref[fielddesc:I2STXMONO]{I2S\_TX\_MONO} = 0, i.e., data is fetched from memory via DMA in both the high and low levels of WS.
    \item \hyperref[fielddesc:I2STXCHANMOD]{I2S\_TX\_CHAN\_MOD} = 2, i.e., mono mode is selected, and the right channel data will be discarded.
    \item \hyperref[fielddesc:I2STXWSIDLEPOL]{I2S\_TX\_WS\_IDLE\_POL} = 1, i.e., both the left channel and right channel transmit the left channel data.
\end{itemize}

Once the configuration is done, assume that the data in memory \textbf{after data format control} is:

\begin{table}[H]
    \centering
    \begin{tabular}{|c|c|c|c|c|c|c|}
    \hline
    Left & Right & Left & Right & ... & Left & Right \\ \hline
    \end{tabular}
\end{table}
\vspace{-0.5em}
\begin{tiplisting}
\vspace{-2em}
\begin{enumerate}
    \item The data above refers to the processed data after data format control instead of the original data.
    \item The ``Left" and ``Right" represent channel data, and their bit widths are channel valid data width. Please refer to Section \ref{sec:i2s-data-control}。
\end{enumerate}
\end{tiplisting}
\vspace{-1em}
Then the channel data is transmitted after channel mode control as follows.

\begin{figure}[H]
    \centering
    \includegraphics[width=0.8\textwidth]{\modulefiles/figures/i2s_chan_mode_20220210.pdf}
    \caption{PDM Channel Control Example}
    \label{fig:i2s-pdm-chan-control-example}
\end{figure}


\section{Receiving Data}\label{sec:i2s-rx-data}


In RX mode, I2S first reads data from peripheral interface, and then stores the data into memory via DMA according to the configured channel mode and data mode.

\subsection{Channel Mode Control}
\chipname{} I2S supports both TDM RX mode and PDM RX mode. Set \hyperref[fielddesc:I2SRXTDMEN]{I2S\_RX\_TDM\_EN} to enable TDM RX mode, or set \hyperref[fielddesc:I2SRXPDMEN]{I2S\_RX\_PDM\_EN} to enable PDM RX mode.

\begin{tiplisting}
\hyperref[fielddesc:I2SRXTDMEN]{I2S\_RX\_TDM\_EN} and \hyperref[fielddesc:I2SRXPDMEN]{I2S\_RX\_PDM\_EN} must not be cleared or set simultaneously.
\end{tiplisting}

\subsubsection{I2S Channel Control in TDM RX Mode}

In TDM RX mode, the total number of RX channels supported is related to the channel valid data width for I2S as follows:

\begin{table}[H]
    \centering
    \caption{The Matching Between Valid Data Width and Number of RX Channel Supported}
    \label{tab:i2s-rx-support-channel}
    \begin{tabular}{|l|r|}
    \hline
    \rowcolor{lightgray}
    \textbf{Channel Valid Data Width} & \textbf{Total Number of Channels Supported} \\ \hline
        32  & 4     \\\hline
        24  & 5     \\\hline
        16  & 8     \\\hline
        8   & 16    \\\hline
    \end{tabular}
\end{table}

In TDM RX mode, I2S supports up to 16 channels to input data. The total number of RX channels in use is controlled by \hyperref[fielddesc:I2SRXTDMTOTCHANNUM]{I2S\_RX\_TDM\_TOT\_CHAN\_NUM}. For example, if  \hyperref[fielddesc:I2SRXTDMTOTCHANNUM]{I2S\_RX\_TDM\_TOT\_CHAN\_NUM} is set to 5, channel 0 \verb+~+ 5 will be used to receive data.

In these RX channels, if \hyperref[fielddesc:I2STXTDMCHAN0EN]{I2S\_RX\_TDM\_CHAN{\textcolor{red}{n}}\_EN} is set to:
\begin{itemize}
    \item 1: this channel data is valid and will be stored into RX FIFO;
    \item 0: this channel data is invalid and will not be stored into RX FIFO.
\end{itemize}
In TDM master mode, WS signal is controlled by \hyperref[fielddesc:I2SRXWSIDLEPOL]{I2S\_RX\_WS\_IDLE\_POL} and \hyperref[fielddesc:I2SRXTDMWSWIDTH]{I2S\_RX\_TDM\_WS\_WIDTH}.
\begin{itemize}
    \item \hyperref[fielddesc:I2SRXWSIDLEPOL]{I2S\_RX\_WS\_IDLE\_POL}: the default level of WS signal;
    \item \hyperref[fielddesc:I2SRXTDMWSWIDTH]{I2S\_RX\_TDM\_WS\_WIDTH}: the cycles the WS default level lasts for when receiving all channel data.
\end{itemize}
\hyperref[fielddesc:I2SRXHALFSAMPLEBITS]{I2S\_RX\_HALF\_SAMPLE\_BITS} x 2 is equal to the BCK cycles in one WS period.

\subsubsection{I2S Channel Control in PDM RX Mode} 
In PDM RX mode, I2S converts the serial data from channels to the data to be entered into memory.

In PDM RX master mode, the default level of WS signal is controlled by \hyperref[fielddesc:I2SRXWSIDLEPOL]{I2S\_RX\_WS\_IDLE\_POL}. WS frequency is half of BCK frequency. The configuration of BCK signal is similar to that of WS signal as described in Section \ref{sec:i2s-tx-rx-clk}. Note, in PDM RX mode, the value of \hyperref[fielddesc:I2SRXHALFSAMPLEBITS]{I2S\_RX\_HALF\_SAMPLE\_BITS} must be same as that of \hyperref[fielddesc:I2SRXBITSMOD]{I2S\_RX\_BITS\_MOD}.


\subsection{Data Format Control}
Data format is controlled in the following phases:
\begin{itemize}
    \item Phase I: serial input data is converted into the data to be saved to RX FIFO;
    \item Phase II: the data is read from RX FIFO and converted according to the input data mode.
\end{itemize}

\subsubsection{Bit Order Control of Channel Data} 
The channel data will be stored as the data to be input in order from high to low. The data bit order in each channel is controlled by \hyperref[fielddesc:I2SRXBITORDER]{I2S\_RX\_BIT\_ORDER}:
\begin{itemize}
\item 0: The bit order of the data to be input is not reversed;
\item 1: The bit order of the data to be input is reversed.
\end{itemize}
At this point, the first phase of data format control is completed.

\subsubsection{Bit Width Control of Channel Storage (Valid) Data}
The storage data width in each channel is controlled by \hyperref[fielddesc:I2SRXBITSMOD]{I2S\_RX\_BITS\_MOD} and \hyperref[fielddesc:I2SRX24FILLEN]{I2S\_RX\_24\_FILL\_EN}. See the table below.

\begin{table}[H]
    \centering
    \caption{Channel Storage Data Width}
    \label{tab:i2s-chan-storage-data-width}
    \begin{tabular}{|p{3cm}|R{3.3cm}|R{4cm}|}
    \hline
    \rowcolor{lightgray}
    \textbf{Channel Storage Data Width} & \textbf{\hyperref[fielddesc:I2SRXBITSMOD]{I2S\_RX\_BITS\_MOD}} & \textbf{\hyperref[fielddesc:I2SRX24FILLEN]{I2S\_RX\_24\_FILL\_EN}}  \\ \hline
    \multirow{2}*{32} & 31 & x \\\cline{2-3}
                 & 23 & 1 \\ \hline
    24 & 23 & 0 \\\hline
    16    & 15 & x \\\hline
    8     & 7 & x \\\hline
    \end{tabular}
\end{table}

\subsubsection{Bit Width Control of Channel RX Data} 
The RX data width in each channel is determined by \hyperref[fielddesc:I2SRXTDMCHANBITS]{I2S\_RX\_TDM\_CHAN\_BITS}.
\begin{itemize}
\item If the storage data width in each channel is smaller than the received (RX) data width, then only the bits within the storage data width is saved into memory. Configure \hyperref[fielddesc:I2SRXLEFTALIGN]{I2S\_RX\_LEFT\_ALIGN} to:
\begin{itemize}
\item 0: only the lower bits of the received data within the storage data width is stored to memory;
\item 1: only the higher bits of the received data within the storage data width is stored to memory.
\end{itemize}
\item If the received data width is smaller than the storage data width in each channel, the higher bits of the received data will be filled with zeros and then the data is saved to memory.
\end{itemize}

\subsubsection{Endian Control of Channel Storage Data} 
The received data is then converted into storage data (to be stored to memory) after some processing, such as discarding extra bits or filling zeros in missing bits. The endian of the storage data is controlled by \hyperref[fielddesc:I2SRXBIGENDIAN]{I2S\_RX\_BIG\_ENDIAN} under various data width. See the table below.
\begin{table}[H]
    \centering
    \caption{Channel Storage Data Endian}
    \label{tab:i2s-chan-storage-data-endian}
    \begin{tabular}{|p{3cm}|p{2.5cm}|p{3.5cm}|R{3.3cm}|}
    \hline
    \rowcolor{lightgray}
    \textbf{Channel Storage Data Width} & \textbf{Original Data} & \textbf{Endian of Processed Data} & \textbf{\hyperref[fielddesc:I2SRXBIGENDIAN]{I2S\_RX\_BIG\_ENDIAN}} \\ \hline
 \multirow{2}*{32} & \multirow{2}*{\{B3, B2, B1, B0\}} & \{B3, B2, B1, B0\} & 0 \\\cline{3-4}
                      && \{B0, B1, B2, B3\} & 1 \\\hline

    \multirow{2}*{24} & \multirow{2}*{\{B2, B1, B0\}} & \{B2, B1, B0\} & 0 \\\cline{3-4}
                      && \{B0, B1, B2\} & 1 \\\hline

    \multirow{2}*{16} & \multirow{2}*{\{B1, B0\}} & \{B1, B0\} & 0 \\\cline{3-4}
                      && \{B0, B1\} & 1 \\\hline

                8     & \{B0\} & \{B0\} & x \\\hline
    \end{tabular}
\end{table}

\subsubsection{A-law/$\mu$-law Compression and Decompression} 

\chipname{} I2S compresses/decompresses the storage data in 32-bit by A-law or by $\mu$-law. By default, zeros are filled into high bits.

Configure \hyperref[fielddesc:I2SRXPCMBYPASS]{I2S\_RX\_PCM\_BYPASS}:
\begin{itemize}
\item 0: compress or decompress the data
\item 1: do not compress or decompress the data
\end{itemize}
Configure \hyperref[fielddesc:I2SRXPCMCONF]{I2S\_RX\_PCM\_CONF}:
\begin{itemize}
\item 0: decompress the data using A-law
\item 1: compress the data using A-law
\item 2: decompress the data using $\mu$-law
\item 3: compress the data using $\mu$-law
\end{itemize}
At this point, the data format control is completed. Data then is stored into memory via DMA.

\section{Software Configuration Process}
\subsection{Configure I2S as TX Mode}\label{sec:i2s-configure-i2s-tx-mode}
Follow the steps below to configure I2S as TX mode via software:

\begin{enumerate}
    \item Configure the clock as described in Section \ref{sec:i2s-tx-rx-clk}.

    \item Configure signal pins according to Table \ref{tab:i2s-signal-descr}.
    
    \item Select the mode needed by configuring \hyperref[fielddesc:I2STXSLAVEMOD]{I2S\_TX\_SLAVE\_MOD}.
    \begin{itemize}
        \item 0: master TX mode
        \item 1: slave TX mode
    \end{itemize}
   
    \item Set needed TX data mode and TX channel mode as described in Section \ref{sec:i2s-tx-data}, and then set \hyperref[fielddesc:I2STXUPDATE]{I2S\_TX\_UPDATE}.

    \item Reset TX unit and TX FIFO as described in Section \ref{sec:i2s-reset}.
    \item Enable corresponding interrupts. See Section \ref{sec:i2s-interrupts}.\label{other:i2s-item-tx-enable-corresponding-interrupts}
    
    \item Configure DMA outlink. 
    
    \item Set \hyperref[fielddesc:I2STXSTOPEN]{I2S\_TX\_STOP\_EN} if needed. For more information, please refer to Section \ref{sec:i2s-master-slave-tx-mode}. 

    \item Start transmitting data:
    \begin{itemize}
        \item In master mode, wait till I2S slave gets ready, then set \hyperref[fielddesc:I2STXSTART]{I2S\_TX\_START} to start transmitting data.
        \item In slave mode, set \hyperref[fielddesc:I2STXSTART]{I2S\_TX\_START}. When the I2S master supplies BCK and WS signals, I2S slave starts transmitting data.
    \end{itemize}

    \item Wait for the interrupt signals set in Step \ref{other:i2s-item-tx-enable-corresponding-interrupts}, or check whether the transfer is completed by querying \hyperref[fielddesc:I2STXIDLE]{I2S\_TX\_IDLE}:
    \begin{itemize}
        \item 0: transmitter is working;
        \item 1: transmitter is in idle state.
    \end{itemize}
    
    \item Clear \hyperref[fielddesc:I2STXSTART]{I2S\_TX\_START} to stop data transfer.
\end{enumerate}
\subsection{Configure I2S as RX Mode}\label{sec:i2s-configure-i2s-rx-mode}
Follow the steps below to configure I2S as RX mode via software:

\begin{enumerate}
    \item Configure the clock as described in Section \ref{sec:i2s-tx-rx-clk}.

    \item Configure signal pins according to Table \ref{tab:i2s-signal-descr}.

    \item Select the mode needed by configuring \hyperref[fielddesc:I2SRXSLAVEMOD]{I2S\_RX\_SLAVE\_MOD}.
    \begin{itemize}
        \item 0: master RX mode
        \item 1: slave RX mode
    \end{itemize}
    
    \item Set needed RX data mode and RX channel mode as described in Section \ref{sec:i2s-rx-data}, and then set \hyperref[fielddesc:I2SRXUPDATE]{I2S\_RX\_UPDATE}.

    \item Reset RX unit and its FIFO according to Section \ref{sec:i2s-reset}.
    \item Enable corresponding interrupts. See Section \ref{sec:i2s-interrupts}.\label{other:i2s-item-rx-enable-corresponding-interrupts}
    
    \item Configure DMA inlink, and set the length of RX data by configuring  \hyperref[regdesc:I2SRXEOFNUMREG]{I2S\_RX\_EOF\_NUM\_REG}.

    \item Start receiving data:
    \begin{itemize}
        \item In master mode, when the slave is ready, set  \hyperref[fielddesc:I2SRXSTART]{I2S\_RX\_START} to start receiving data.
        \item In slave mode, set  \hyperref[fielddesc:I2SRXSTART]{I2S\_RX\_START} to start receiving data when get BCK and WS signals from the master.
    \end{itemize}
    \item The received data is then stored to the specified address of \chipname{} memory according the configuration of DMA. Then the corresponding interrupt set in step \ref{other:i2s-item-rx-enable-corresponding-interrupts} is generated.
\end{enumerate}


\section{I2S Interrupts} \label{sec:i2s-interrupts}

\begin{itemize}
\item \label{int:I2STXHUNGINT}I2S\_TX\_HUNG\_INT: triggered when transmitting data is timed out. For example, if module is configured as TX slave mode, but the master does not provide BCK or WS signal for a long time (specified in \hyperref[regdesc:I2SLCHUNGCONFREG]{I2S\_LC\_HUNG\_CONF\_REG}), then this interrupt will be triggered.
\item \label{int:I2SRXHUNGINT}I2S\_RX\_HUNG\_INT: triggered when receiving data is timed out. For example, if I2S module is configured as RX slave mode, but the master does not send data for a long time (specified in \hyperref[regdesc:I2SLCHUNGCONFREG]{I2S\_LC\_HUNG\_CONF\_REG}), then this interrupt will be triggered.
\item \label{int:I2STXDONEINT}I2S\_TX\_DONE\_INT: triggered when transmitting data is completed.
\item \label{int:I2SRXDONEINT}I2S\_RX\_DONE\_INT: triggered when receiving data is completed.
\end{itemize}

\section{Event Task Matrix Feature}\label{sec:i2s-etm-func}

\chipname{} I2S supports the Event Task Matrix (ETM) function, which allows I2S's ETM tasks to be triggered by any peripherals' ETM events, or I2S's ETM events to trigger any peripherals' ETM tasks. This section introduces the ETM tasks and events related to I2S. For more information, please refer to Chapter \ref{mod:evntaskmatrix} \nameref{mod:evntaskmatrix}.

I2S can receive the following ETM tasks:

\begin{itemize}
    \item \label{task:I2STASKSTARTTX}I2S\_TASK\_START\_TX: Enables I2S TX for data transfer.
    \item \label{task:I2STASKSTARTRX}I2S\_TASK\_START\_RX: Enables I2S RX for data transfer.
    \item \label{task:I2STASKSTOPTX}I2S\_TASK\_STOP\_TX: Stops I2S TX data transfer.
    \item \label{task:I2STASKSTOPRX}I2S\_TASK\_STOP\_RX: Stops I2S RX data transfer.
\end{itemize}

I2S can generate the following ETM events:

\begin{itemize}
    \item \label{evt:I2SEVTTXDONE}I2S\_EVT\_TX\_DONE: Indicates that I2S TX has completed data transmission.
    \item \label{evt:I2SEVTRXDONE}I2S\_EVT\_RX\_DONE: Indicates that I2S RX has completed data receiving.
    \item \label{evt:I2SEVTXWORDSSENT}I2S\_EVT\_X\_WORDS\_SENT: Indicates that the word number sent by I2S TX is equal to or larger than the value set by \hyperref[fielddesc:I2SETMTXSENDWORDNUM]{I2S\_ETM\_TX\_SEND\_WORD\_NUM}.
    \item \label{evt:I2SEVTXWORDSRECEIVED}I2S\_EVT\_X\_WORDS\_RECEIVED: Indicates that the word number received by I2S RX is equal to or larger than the value set by \hyperref[fielddesc:I2SETMRXRECEIVEWORDNUM]{I2S\_ETM\_RX\_RECEIVE\_WORD\_NUM}.
\end{itemize}

In practical applications, I2S's ETM events can trigger its own ETM tasks. For example, the I2S\_EVT\_X\_WORDS\_SENT event can trigger the I2S\_TASK\_STOP\_TX task, and in this way stop the I2S operation through ETM.



% >>> If your module has no registers, delete sample text
%       for Sections Register Summary and Registers below <<<

%%%%%%%%%%%%%%%%%%%%%%%%
%%% Register Summary %%%
%%%%%%%%%%%%%%%%%%%%%%%%

% >>> Update the sample text below and insert
%     the info on register summary into the subfile <<<

\clearpage
\hypertarget{i2s-reg-summ}{}
\section{Register Summary}
\label{sec:i2s-reg-summ}

The addresses in this section are relative to I2S Controller base address provided in Table \ref{tab:sysmem-base-address} in Chapter \ref{mod:sysmem} \textit{\nameref{mod:sysmem}}.

The abbreviations given in Column \textbf{Access} are explained in Section \hyperref[glossary-access-types]{\textit{Access Types for Registers}}.

\subfile{\modulefiles/33-I2S-reg-summ__EN}


%%%%%%%%%%%%%%%%%
%%% Registers %%%
%%%%%%%%%%%%%%%%%

% >>> Update the sample text below and insert
%      the info on registers into the subfile <<<

\clearpage
\section{Registers}

The addresses in this section are relative to I2S Controller base address provided in Table \ref{tab:sysmem-base-address} in Chapter \ref{mod:sysmem} \textit{\nameref{mod:sysmem}}.

\subfile{\modulefiles/33-I2S-reg__EN}


\end{document}
